% 03-resolution.tex — Résolution numérique (XPBD à petits pas, solveur direct).

\section{Résolution numérique}

Une corde inextensible modélisée par des ressorts très raides demanderait un
pas de temps minuscule. La dynamique basée sur les positions
\cite{muller2007pbd, bender2014survey} contourne ce problème : à chaque pas, on
prédit les positions sans tenir compte des liaisons, puis on les
\emph{projette} sur les liaisons, et l'on déduit enfin les vitesses des
déplacements. Sa forme étendue XPBD \cite{macklin2016xpbd} associe à chaque
liaison un multiplicateur de Lagrange \(\lambda_i\), d'où l'on tire la force :
la tension d'un maillon vaut \(T_i = -\lambda_i / \Delta t^2\). Un pas
\(\Delta t\) comporte quatre étapes :
\begin{enumerate}
  \item prédiction : \(\vect{x}^* = \vect{x} + \vect{v} \Delta t + \vect{g}_{\text{eff}} \Delta t^2\) pour chaque point libre (\(\vect{g}_{\text{eff}} = \vect{g}\), sauf \(3\vect{g}/2\) pour l'extrémité de la barre) ;
  \item projection de toutes les liaisons de la chaîne en une fois (ci-dessous) ;
  \item projection des gouttes sur les parois du seau et entre elles ;
  \item vitesses : \(\vect{v} = (\vect{x}^* - \vect{x}) / \Delta t\), puis frottement aux articulations.
\end{enumerate}

\paragraph{Une résolution directe plutôt qu'itérative.} La méthode PBD
classique projette les liaisons une par une (Gauss--Seidel). Pour un seau
lourd au bout de nœuds de corde légers, une passe ne suffit pas : la corde
s'allonge. Avec un seau de \SI{5}{kg} au bout d'une corde de \SI{30}{cm}, une
passe de Gauss--Seidel laisse la corde s'allonger jusqu'à \SI{1,04}{m} en
\SI{0,38}{s}, plus de trois fois sa longueur, alors que l'écart reste inférieur
à \SI{e-5}{mm} avec la résolution directe (données de
\texttt{outils/donnees.mjs}). La linéarisation des
contraintes \(C_i = |\vect{x}_{i+1} - \vect{x}_i| - \ell_i\) donne le système
\(J M^{-1} J^{\mathsf{T}} \Delta\vect{\lambda} = -\vect{C}\). Comme la liaison
\(i\) ne touche que les points \(i\) et \(i+1\), cette matrice est
tridiagonale : elle se résout en \(O(n)\) opérations par l'algorithme de Thomas
\cite{press2007}, comme le font les solveurs directs de tiges
\cite{deul2018direct}. Les maillons de corde sont unilatéraux : un maillon
détendu (\(C_i < 0\)) ou qui devrait pousser est retiré du système (ensemble
actif), et le système est résolu à nouveau.

\paragraph{Petits pas et stabilité.} Plutôt que plusieurs itérations par pas,
la simulation fait beaucoup de pas courts \cite{macklin2019smallsteps} : au
moins \num{1200} par seconde pendant l'élan et \num{2400} pendant le lancer et
les tours libres. Le pas doit aussi suivre les ondes transverses de la corde,
de célérité \(c = \sqrt{T/\mu}\) : la fréquence des pas vaut au moins
\(1{,}3\, c/\ell\), ce qui compte pour une corde courte sous un seau lourd. La
tension affichée est celle du maillon le plus chargé, filtrée sur \SI{20}{ms}.
La projection dissipe un peu d'énergie, de 2 à \SI{8}{\percent} en deux tours.
